
A four-hypothesis validation chain was designed to test whether incremental SMT verification would reduce verification time for LLM code repair in typed Python with Pydantic. This section describes the hypothesis decomposition, experimental protocol, and pipeline architecture. While validation was blocked by infrastructure failures, the methodology and implementation artifacts remain sound.

\subsubsection{Hypothesis Decomposition}

\textbf{Main Hypothesis (H-IncrementalSMT-v1)}: Under iterative LLM code repair workflows in typed Python with Pydantic, if incremental SMT verification is used (re-verifying only modified functions and dependencies), then verification time reduces compared to batch re-verification, because unchanged code portions skip redundant constraint checking.

This decomposes into four sub-hypotheses:

\begin{itemize}
\item \textbf{h-e1 (EXISTENCE)}: Static analyzers can extract SMT constraints from LLM-generated typed code. Gate: $\geq 90$\% extraction success rate. If fail: pivot to neural constraint extraction. \textbf{Status: FAIL (0.0\%, infrastructure block).}
\end{itemize}

\begin{itemize}
\item \textbf{h-m1 (MECHANISM)}: Pyre extracts type contracts from Pydantic annotations. Gate: MUST\_WORK. \textbf{Status: NOT\_STARTED (prerequisite h-e1 failed).}
\end{itemize}

\begin{itemize}
\item \textbf{h-m2 (MECHANISM)}: Dependency analysis computes transitive closure for incremental re-verification. Gate: SHOULD\_WORK. \textbf{Status: NOT\_STARTED.}
\end{itemize}

\begin{itemize}
\item \textbf{h-m3 (MECHANISM)}: Incremental SMT achieves speedup vs batch. Gate: MUST\_WORK. \textbf{Status: NOT\_STARTED.}
\end{itemize}

The dependency chain is sequential: h-e1 $\rightarrow$ h-m1 $\rightarrow$ h-m2 $\rightarrow$ h-m3. Gate failure at h-e1 blocked all downstream hypotheses.

\subsubsection{Dataset Extension: HumanEval + Pydantic}

HumanEval lacks type annotations required for SMT constraint extraction. A template-based transformation pipeline was designed:

\textbf{Step 1}: Parse function signature and docstring from each HumanEval problem.

\textbf{Step 2}: Generate Pydantic BaseModel schema for inputs and outputs:
\begin{verbatim}
class Input(BaseModel):
    param1: Type1
    param2: Type2
    
    @validator('param1')
    def validate_param1(cls, v):
        assert v is not None, 'param1 must not be None'
        return v
\end{verbatim}

\textbf{Step 3}: Add @validator decorators encoding preconditions extracted from docstrings.

\textbf{Step 4}: Preserve original test suite for functional correctness validation.

This process was applied to HumanEval problems 0--99, generating 100 Pydantic-annotated prompts. All 100 prompts were created successfully, demonstrating feasibility of mechanical type annotation generation.

\textbf{Design Rationale}: Pydantic validators provide runtime-checkable contracts that map naturally to SMT predicates. Mypy type hints support static checking but lack explicit constraint semantics. HumanEval was chosen because it provides existing test suites and serves as a standard baseline for code generation research. Template-based transformation ensures consistent annotation quality and scales to large benchmarks, though templates encode only simple constraint patterns (non-null, range checks). Complex preconditions require semantic understanding beyond pattern matching.

\subsubsection{Constraint Extraction Pipeline}

The pipeline consists of five sequential modules:

\begin{enumerate}
\item \textbf{DatasetExtender}: Load HumanEval, generate Pydantic templates (150 LOC).
\item \textbf{LLMGenerator}: Generate typed Python code via Claude Sonnet 3.5 API (82 LOC).
\item \textbf{PyreExtractor}: Parse AST, extract SMT constraints from type annotations (128 LOC).
\item \textbf{Z3Validator}: Validate constraint quality via SMT satisfiability checking (132 LOC).
\item \textbf{MetricsEngine}: Compute extraction/quality rates, generate visualizations (156 LOC).
\end{enumerate}

Modular design allowed partial validation despite LLM generation failure. DatasetExtender was validated on 100 prompts (success). LLMGenerator failed with HTTP 401 errors (invalid API key). PyreExtractor was validated on error files (correctly found 0 constraints). Z3Validator was validated on error files (correctly identified 0 quality constraints). MetricsEngine correctly computed 0.0\% extraction rate and generated 4 figures.

Pipeline validated on negative cases only. Pipeline separation (extend $\rightarrow$ generate $\rightarrow$ extract) prevented cascading failures---dataset extension succeeded independently of generation failure.

\subsubsection{Incremental SMT Strategy}

\textbf{Batch Verification (Baseline)}: On each repair iteration, re-verify entire program by re-extracting all constraints and solving via Z3.

\textbf{Incremental Verification (Proposed)}:
\begin{enumerate}
\item Extract constraints from initial program, solve via Z3 push/pop contexts.
\item On modification, run dependency analysis to compute invalidation cone (functions reachable from modified code).
\item Re-verify only invalidated constraints; reuse cached results for unchanged code.
\end{enumerate}

\textbf{Success Criteria}: Not validated due to infrastructure failure.

Conservative dependency analysis was designed for soundness: over-approximate dependency cone to prioritize soundness. False positives (re-verifying more than necessary) reduce speedup but maintain correctness. False negatives (missing dependencies) cause unsound verification. Z3 incremental mode uses push/pop contexts to cache verification state across iterations. The approach assumes localized repairs (based on CURE/CoCoNut literature showing 80\% of human repairs touch 1--3 lines), but this assumption remains unverified for LLM code.

\subsubsection{Validation Protocol}

\textbf{h-e1 Protocol}:
\begin{enumerate}
\item Generate 100 typed Python programs from Pydantic-annotated HumanEval prompts using Claude Sonnet 3.5.
\item Extract SMT constraints via Pyre static analyzer.
\item Measure extraction success rate (\% of programs with usable constraints).
\item Compare against 90\% threshold.
\end{enumerate}

\textbf{Actual Result}: Step 1 failed with 100\% API authentication errors. extraction\_rate = 0.0\% (infrastructure block, not extraction failure).

h-m1, h-m2, h-m3 protocols were not executed due to prerequisite h-e1 failure.

\subsubsection{Failure Mode Analysis}

The h-e1 validation failed with 0.0\% extraction rate due to 100\% API authentication failure (48 generation attempts returned HTTP 401 ''API key is invalid''). Root cause: invalid Anthropic API key in experiment configuration. The retry mechanism (3 attempts with exponential backoff) correctly exhausted retries but lacked logic to fast-fail on non-retryable auth errors.

The falsifier for h-e1 is ''static analyzer fails to extract usable constraints (<90\% success).'' Infrastructure failures prevent testing the falsifier. The hypothesis remains theoretically plausible pending retry with valid API access. The pipeline is complete and validated on negative cases. Retry requires only fixing the API key, not redesigning the experiment.

\subsubsection{Reproducibility}

\textbf{Environment}: Python 3.13, dependencies: \texttt{datasets}, \texttt{anthropic}, \texttt{pydantic}, \texttt{z3-solver}, \texttt{matplotlib}.

\textbf{Data}: HumanEval (openai\_humaneval, problems 0--99), extended prompts in \texttt{src/data/humaneval\textbackslash \_pydantic/}.

\textbf{Code}: 700 LOC total, modular architecture.

\textbf{Retry Instructions}: Set valid \texttt{ANTHROPIC\textbackslash \_API\textbackslash \_KEY}, run \texttt{python main.py}. Expected outcome (if hypothesis valid): extraction\_rate $\geq 90$\%, gate PASS.
